%Do not modify this content without consulting with its authors!!!
%Adam and Yen
\subsection{Challenges and Limitations}
While generative AI offers transformative potential, its deployment introduces a complex landscape of challenges spanning technical, ethical, social, and governance dimensions. Our thematic analysis reveals four primary categories of concerns that demand careful attention to ensure responsible and effective implementation of these technologies.

\subsubsection{Societal, Ethical and Fairness Concerns [C1]}
\paragraph{Biases and Discrimination}
A predominant ethical challenge, identified consistently across the literature, is the propensity of GenAI models to perpetuate and amplify societal biases. This concern transcends domain boundaries, manifesting across general management information systems \parencite{R01, R02, R04, R07, R08, R10, S09}, healthcare applications \parencite{R09, S01, S11, S16}, and educational contexts \parencite{S14}. The problem's origin lies in the models' training data, which often encapsulates historical and systemic prejudices. As \citeauthor{R02} explain,\enquote{the training data of LLMs may contain biases from various sources reflecting racial, gender, and other discriminant judgments in humans and society. Trained on these data, LLMs may inherit and amplify such biases, causing the decisions to be unfair for some social groups, communities, or societies} \parencite{R02}. 

The ramifications of this inherited bias are severe, leading to discriminatory outputs and representational harms that disproportionately affect marginalized groups \parencite{R02, S09, S11}. This is compounded by what one study terms ``exclusionary norms'', where models trained on data from affluent regions neglect global diversity, thereby reflecting the \enquote{practices of the wealthiest communities and countries} and fostering cultural insensitivity \parencite{S09}. The tangible risks of such biases are particularly acute in high-stakes applications. In healthcare, they can manifest as clinically inappropriate recommendations stemming from a failure to grasp linguistic or cultural nuances \parencite{S01}. In organizational settings, biased algorithms can unfairly influence critical decisions like hiring and firing \parencite{R09}, while in education, they risk reinforcing discriminatory worldviews among learners \parencite{S14}.

\paragraph{Fairness Implementation Challenges}
Addressing bias is not merely a technical problem of detection but a profound normative challenge of implementation, fraught with both conceptual and practical barriers. A primary conceptual hurdle is the absence of a universal definition of fairness, as what is considered equitable is deeply embedded in cultural, legal, and social norms \parencite{R01}. This definitional ambiguity becomes particularly salient in content moderation, where, as one study highlights, \enquote{there is no universally accepted definition of what qualifies as hate speech or toxic speech} \parencite{S09}. Without clear, agreed-upon criteria for what constitutes harmful content, creating globally consistent and fair moderation policies becomes exceptionally difficult.

Compounding this normative challenge are practical tensions, as the goals of fairness and equity often conflict with organizational objectives such as profitability and operational efficiency \parencite{R01, R08}. These intertwined conceptual and practical obstacles mean that even when biases are identified, rectifying them in a consistent and meaningful way remains a formidable task.

\paragraph{Potential for Misuse and Harmful Content}
GenAI systems present unprecedented capabilities for generating harmful, manipulative, and malicious content at scale. The scope of potential exploitation is expansive. \citeauthor{S09} provide a comprehensive catalog of deliberate misuse scenarios, ranging from the generation of \enquote{malevolent material, including spam, fraudulent reviews, or even cyberattacks on a large scale} to \enquote{creating deceptive phishing emails and malicious code} \parencite*{S09}. Deepfake technology represents a particularly insidious vector, enabling sophisticated identity fraud and deception, while GenAI's capacity for emotional manipulation introduces novel forms of psychological harm \parencite{R10}.

The propagation of misinformation and disinformation represents another pressing concern \parencite{R03, R08}. GenAI models \enquote{risk blurring the line between fact and fiction, as they can rapidly disseminate false or misleading information, fake news, and malicious content, making it difficult for users to discern truth from fantasy} \parencite{S09}. The consequences of such information pollution vary dramatically by domain. In healthcare contexts, misinformation can directly endanger patient safety through inaccurate medical guidance \parencite{S01}, while in political spheres, GenAI \enquote{can be exploited for manipulative purposes, such as the generation of propaganda or misinformation, thereby influencing public opinion and potentially harming, for example, the electoral process or other fraud and scams} \parencite{S09}. Educational settings reveal different manifestations of misuse, including examination fraud and plagiarism that undermine academic integrity \parencite{S12, S14}.

Underlying these diverse threats is the observation that GenAI tools lack a moral compass, operating without the \enquote{intuition, plausibility, and temporal relevance} that guide human judgment \parencite{R06}. This technological amorality, combined with the potential for widespread misuse, erodes societal trust in institutions and information ecosystems \parencite{S09}. Consequently, it raises urgent questions about accountability and the difficulty of moderating AI-generated content at scale \parencite{S01, S09, R07, S17}.

\paragraph{Human Autonomy and Skill Devaluation}
A recurring theme in the literature is the concern that over-reliance on GenAI may erode essential human skills and autonomy \parencite{S01, S08, R09, S17}. The primary mechanism for this is the development of a \enquote{human automation bias,} where users accept AI-generated answers without critical assessment, leading to a dependency that \enquote{risks eroding skills such as creativity and critical thinking} \parencite{S09}. This risk is particularly pronounced in educational settings, where the ease of content generation is seen as a threat to the development of students' independent reasoning, perseverance, and resilience \parencite{S12, S14}. Beyond cognitive decline, the literature also points to the degradation of social bonds, as increased automation may reduce meaningful human interaction and lead to more profound harms, including the \enquote{loss of autonomy and dignity, dehumanization, social isolation, and addiction} \parencite{R08, S11}.

This erosion of autonomy is compounded by the human tendency to anthropomorphize conversational AI, which introduces distinct ``psychological vulnerabilities'' \parencite{S09}. When users perceive AI systems as human-like, they become more susceptible to developing inappropriate dependencies and may be more likely to disclose sensitive personal information \parencite{S09}. By fostering a false sense of relationship, anthropomorphism can deepen the very risks of skill erosion and social isolation previously discussed, blurring the lines between tool and companion in ways that may undermine user agency.

\paragraph{Environmental Sustainability}
The environmental costs of GenAI constitute an often-overlooked yet critical dimension of ethical deployment. Training and operating large-scale generative models impose substantial resource demands, resulting in significant energy consumption and carbon footprints \parencite{R01, R04, S07, S10, R10, S15}. As \citeauthor{S09} emphasize, \enquote{the energy demands for training and operating these models contribute to resource depletion and pollution, leave a significant carbon footprint, and have high computation costs} \parencite*{S09}. This challenge is intrinsically linked to what \citeauthor{R02} describe as \enquote{the blessing and curse of the scaling law}—the empirical observation that model performance improves with increases in size, training data, and computational power \parencite*{R02}. This dynamic creates perverse incentives that encourage ever-larger and more resource-intensive architectures, establishing a tension between performance optimization and environmental sustainability.

Compounding these concerns is the research community's limited visibility into the full scope of environmental impacts. \citeauthor{S09} observe that \enquote{many environmental factors related to the operation of LLMs that are in widespread use are currently unknown} \parencite*{S09}, implying that documented concerns may represent merely a fraction of the true ecological cost. This opacity raises fundamental questions about the sustainability of current GenAI development trajectories and highlights an ethical imperative to systematically account for environmental impacts alongside performance metrics \parencite{S15}.

\subsubsection{Reliability and Accuracy [C2]}
%A fundamental challenge pervading the literature is the questionable reliability and accuracy of GenAI outputs, which presents a significant barrier to its adoption in mission-critical applications across all reviewed domains \parencite{R04, S05, S16, S17}. This concern manifests in several distinct but interrelated forms, ranging from the generation of outright falsehoods to more subtle failures in contextual understanding and performance consistency.

\paragraph{Hallucinations and Factual Inaccuracy}
The most pervasive technical challenge undermining GenAI reliability is the phenomenon of ``hallucination''—the generation of outputs that appear plausible but are fundamentally incorrect \parencite{R02,R03,R08,R10,S09}. This issue manifests consistently across all application domains examined, from healthcare \parencite{S01, S02, S03, S05, S10, S11, S16, S17} and education \parencite{S14} to general information systems \parencite{R02, R03, R04, R05, R06, R07, R08}.

The consequences of such inaccuracies vary dramatically by context, with healthcare applications facing the greatest risks. \citeauthor{S02} warn that \enquote{inaccurate or misleading information in healthcare can have severe consequences, including misdiagnoses, improper treatments, and potential harm to patients' well-being and safety} \parencite*{S02}. Empirical evidence substantiates this concern: \citeauthor{S05} documented \enquote{the presence of hallucinations or fictitious information in three of nine studies utilizing ChatGPT} \parencite*{S05}, raising fundamental questions about safe clinical implementation. The fabrication problem extends beyond medical contexts to academic and professional settings, where models have been observed generating fictitious references \parencite{S14}. This phenomenon of hallucination renders human oversight and validation indispensable for any responsible application of GenAI \parencite{R02, R05}.

\paragraph{Outdated or Limited Knowledge}
A structural constraint on GenAI reliability stems from the temporal constraints of their training data, which has a fixed knowledge cut-off date \parencite{R04, R06, S14}. This temporal freezing creates an expanding knowledge gap as models age, rendering them increasingly obsolete for queries requiring current information. This inherent limitation is compounded by the fact that models \enquote{may not remember everything that they saw during training} \parencite{R04}, leading to incomplete or sparse knowledge—a particular challenge when domain-specific expertise is required \parencite{S07, S10, S11}.

These knowledge gaps manifest differently across domains: in healthcare as an inability to answer complex medical questions \parencite{S11}, in research as dependency on limited or unrepresentative datasets \parencite{S07}, and in general applications as an absence of temporal awareness that undermines contextual relevance \parencite{R06}. Furthermore, the training process itself amplifies these limitations, as developers \enquote{often tend to rely on second-hand information from official organizations, which has a certain degree of authority but is often lagging behind} \parencite{S11}. This creates a cascade of temporal delays—from the original data collection to model training to user deployment—each stage introducing additional staleness into the system.

\paragraph{Limited Contextual Awareness}
Beyond factual accuracy, a more subtle yet critical challenge is GenAI's limited contextual awareness—an inability to interpret situational, cultural, or domain-specific nuances \parencite{S01, R04, R05, S07, S09, S11}. This deficiency can render outputs technically correct yet practically useless, inappropriate, or even harmful. The risk is particularly acute in healthcare, where this limitation manifests as a failure to navigate linguistic and cultural subtleties \parencite{S01} or provide genuinely personalized health information \parencite{S11}.

\paragraph{Performance Inconsistency and Drift}
The reliability of GenAI is further complicated by its dynamic and often unpredictable nature. The literature highlights concerns about performance ``drift'', described as the unexpected deterioration of model performance over time \parencite{R04}, as well as inconsistent reliability across different clinical scenarios \parencite{S05}. This instability also appears at a technical level, with some models exhibiting instability during training \parencite{S07} or vulnerability to \enquote{semantic perturbations, whereby input with different syntax but similar meaning to the training data leads to errors} \parencite{S09}, and susceptibility to system crashes \parencite{S11}.

The combination of performance drift and inherent instability creates a particularly challenging scenario for deployment, as initial testing may not reveal latent failure modes that emerge over time or under specific conditions. This unpredictability undermines trust and necessitates constant vigilance, transforming what are marketed as autonomous systems into tools requiring continuous human oversight and validation \parencite{R04}.

\paragraph{Lack of Evidence and Validation}
Underpinning all other reliability concerns is a critical meta-challenge: the lack of rigorous, evidence-based validation of GenAI systems in real-world applications. This issue is particularly pronounced in specialized domains like healthcare, where there is a \enquote{scarcity of evidence-based medical research concerning the application of LLMs in healthcare settings} \parencite{S01}. The validation gap extends across multiple dimensions: absence of external validation studies, a lack of comprehensive and domain-specific evaluation metrics, and the immaturity of assessment methods and theoretical frameworks \parencite{S01, S07, S08, S18}.

Without established benchmarks and empirical evidence, practitioners find it difficult to delineate when these tools are productive and when they are not \parencite{R06}. Compounding this challenge is a lack of reproducibility, as it may not be possible to reliably replicate tool responses through prompt engineering \parencite{R06}. These validation deficiencies not only hinder safe and effective integration \parencite{S01, S07} but also underscore the urgent need for rigorous, domain-specific evaluation frameworks before widespread deployment \parencite{S07, S18}.

\subsubsection{Transparency and Explainability [C3]}
\paragraph{Model Opacity and the Black-Box Problem}
The inherent opacity of GenAI models, frequently described as a ``black box'' problem, stems from the difficulty in interpreting how their complex transformer-based architectures arrive at specific outputs \parencite{R02, S09, R10}. This lack of transparency is a pervasive concern that obstructs the ability to audit decision-making processes \parencite{R02, R06}, assess model limitations \parencite{R08, S14}, and ensure user control \parencite{R09}, eroding trust across all reviewed domains \parencite{S08, R08, S14, S16}. Consequently, stakeholders are confronted with systems of immense size and ``opaque behaviors'' \parencite{S09} whose operational logic remains inscrutable, hindering both accountability and safe adoption.

\paragraph{Explainability Shortcomings}
Compounding the problem of model opacity are significant deficiencies in the tools and methodologies for explaining GenAI systems \parencite{S06, S07, S09, S18}. Existing Explainable AI (XAI) techniques are described as \enquote{still far from optimal,} with a general consensus that available tools are marked by their ``immaturity'' \parencite{S18}. This deficiency makes it difficult to \enquote{interpret and explain the rationale behind the decision-making process of a model} \parencite{R02} or its \enquote{non-interpretable learned representations} \parencite{S07}. The challenge is further amplified by the technology's scale, as \enquote{XAI for GenAI faces significant challenges due to the increasing complexity and societal reach of these models} \parencite{S18}. Ultimately, without robust explainability, it is nearly impossible to debug, validate, or ethically govern GenAI systems, leaving a critical gap between their advanced capabilities and the human capacity to responsibly manage them.

\subsubsection{Privacy, Security, Governance and Legal [C4]}
\paragraph{Data Privacy Violations}
A pervasive concern across all reviewed domains is the risk of unintended leakage or non-consensual exposure of sensitive, personal, or proprietary information \parencite{S01, S02, R08, R09, S08, S11, S15, S16, S17}. This risk originates from multiple sources, beginning at the model level where systems demonstrate \enquote{a tendency to memorize and reproduce personally identifiable information} from their training data \parencite{S09}.

These inherent risks are then amplified by systemic vulnerabilities in deployment. \citeauthor{S09} report that \enquote{owing to system glitches in ChatGPT, the chat logs of certain users have become accessible to others,} affecting both individuals and organizations. Similarly, in enterprise contexts, organizations hesitate to use public AI tools because the prompts they submit can reveal sensitive information \parencite{R03}. This risk of disclosure is further compounded by a behavioral dimension, as users are more likely to reveal private information when they anthropomorphize the technology and \enquote{treat models as if they are human} \parencite{S09}. Collectively, these ``privacy hazards'' \parencite{S09} create significant challenges for compliance with regulations like GDPR and HIPAA \parencite{S09, S16}.

\paragraph{Security Vulnerabilities}
GenAI systems are also susceptible to security vulnerabilities that expose them to malicious exploitation \parencite{R03, S08, S10, S16, S17}. \citeauthor{S09} document a range of adversarial attack vectors designed to compromise system integrity, including susceptibility to ``jailbreaking'', where prompts are used to bypass safety measures; ``prompt injection'' to cause malfunctions; and ``data poisoning attacks'' to corrupt the model's knowledge base \parencite{S09}. Such security breaches create pathways for unauthorized access, data theft, and other intentional misuses that threaten both personal and organizational security \parencite{S08, S09}.

\paragraph{Regulatory and Legal Gaps}
The rapid proliferation of GenAI has created a significant governance vacuum, as existing legal frameworks are ill-equipped to manage the technology. The literature consistently notes that \enquote{current laws and regulations have become inadequate to account for new phenomena brought about by [GenAI]} \parencite{R07}, with the technology being adopted far faster than it can be theorized or governed \parencite{S13}. This creates a profound lack of globally agreed-upon standards for ethical and safe deployment \parencite{R09}. This regulatory gap poses a direct challenge for organizations attempting to ensure legal compliance \parencite{S08} and for regulatory bodies tasked with developing policies to protect the public, particularly in high-stakes domains like healthcare \parencite{S16}.

\paragraph{Accountability and Liability Ambiguity}
The legal and regulatory vacuum directly contributes to a critical ambiguity regarding accountability and liability. When a GenAI system causes harm through an error or biased output, there is no clear framework for assigning responsibility, leaving a notable ``ambiguity over liability'' \parencite{S17}. This uncertainty extends to defining moral responsibility for model outputs \parencite{S09} and is exemplified by the practical ``warranty problem'', where model suppliers refuse to provide performance guarantees, forcing adopters to shoulder the operational and legal risks \parencite{R03}.

\paragraph{Copyright and Ownership Disputes}
Finally, GenAI systems raise novel challenges to established notions of intellectual property, creating unresolved disputes over both model inputs and outputs \parencite{R04, R08, R10, S10, S12, S14}. On the input side, models are trained on vast datasets that may include copyrighted material without permission, potentially violating existing rights \parencite{R04, S09}. On the output side, the \enquote{distinction between original and AI-generated content is blurred} \parencite{S09}, leading to profound \enquote{doubts about who is a legal owner of GenAI generated contents} \parencite{R04}. This ambiguity fuels practical concerns about plagiarism and academic integrity \parencite{S10,S12} and complicates questions of ownership and control in commercial contexts \parencite{S09}.
